% ============================================================
% PARTE V - LA DISSONANZA: NOTE ESTRANEE E RITARDI
% CAPITOLO 18 - TEORIA GENERALE DELLA DISSONANZA
% ============================================================
\part{La dissonanza: note estranee e ritardi}

\chapter{Teoria generale della dissonanza}
\label{chap:teoria_dissonanza}

Finora gli accordi erano consonanti. La dissonanza è l'elemento dinamico che genera tensione, movimento e bisogno di risoluzione. Questo capitolo fissa il protocollo con cui una nota dissonante entra e si risolve; la stessa logica regge le note estranee, i ritardi e gli accordi di settima.

\section{Il ciclo della dissonanza}
\label{sec:ciclo_dissonanza}

Nello stile severo e nella teoria classica la gestione di una nota dissonante segue un protocollo in tre fasi.

\begin{regola}[Trattamento della dissonanza]
	\label{reg:trattamento_dissonanza}
	Ogni nota dissonante attraversa tre momenti distinti:
	\begin{enumerate}
		\item \textbf{Preparazione:} la nota che diventerà dissonante è presente nell'accordo precedente come \textbf{consonanza}, di solito mantenuta nella stessa voce. La preparazione giustifica l'ingresso della tensione.
		\item \textbf{Percussione:} la nota è suonata insieme alle altre note dell'accordo e diventa effettivamente \textbf{dissonante} rispetto alla fondamentale o ad altre voci.
		\item \textbf{Risoluzione:} la nota dissonante si muove per \textbf{grado congiunto discendente} verso una nota consonante nell'accordo successivo.
	\end{enumerate}
\end{regola}

\begin{esempio}[Preparazione e risoluzione]
	\label{ex:ciclo_dissonanza}
	Qui compare un accordo di settima, che si studia dal Capitolo \ref{chap:settima_dominante}; interessa soltanto il percorso della nota dissonante. Nel collegamento I-\sigla{II}{7}{}-V il Do del soprano è consonante nel primo accordo (ottava), diventa dissonante nel secondo (settima del Re) e risolve scendendo al Si (terza del Sol).
	\musica{es_cap08_01}
\end{esempio}

\section{Evoluzione del concetto di preparazione}
\label{sec:evoluzione_preparazione}

La \textbf{risoluzione discendente} resta una regola quasi inviolabile dell'armonia tonale. La \textbf{preparazione}, invece, ha perso importanza con l'evoluzione del linguaggio.

Alcune dissonanze, per la forte direzionalità e la familiarità (come la settima di dominante), hanno cominciato a essere introdotte «di petto», senza preparazione. Anche allora la nota dissonante conserva la sua natura e la sua esigenza di risoluzione.

\begin{regola}[Dissonanze non preparate]
	\label{reg:dissonanze_non_preparate}
	Gli accordi di dominante (\sigla{V}{7}{}, \sigla{VII}{7}{} e derivati) ammettono spesso la percussione della dissonanza senza preparazione. Gli accordi di settima delle altre specie (\sigla{II}{7}{}, \sigla{IV}{7}{} e così via) richiedono quasi sempre una preparazione rigorosa nello stile scolastico.
\end{regola}

Il protocollo di questo capitolo si applica a:
\begin{itemize}
	\item le note estranee (Capitolo \ref{chap:noteestranee}) e i ritardi (Capitolo \ref{chap:ritardi});
	\item la settima di dominante (Capitolo \ref{chap:settima_dominante}), dove la risoluzione è il motore della cadenza;
	\item il \sigla{VII}{7}{} e il \sigla{II}{7}{} (Capitolo \ref{chap:vii7_ii7}), dove la preparazione conta per la fluidità del discorso;
	\item tutte le altre specie di settima (Capitolo \ref{chap:sette_specie}), che si leggono come estensioni dello stesso principio.
\end{itemize}
